\documentclass[12pt]{article}
\usepackage[utf8]{inputenc}
\usepackage[T1]{fontenc}
\usepackage{hyperref}
\title{Akei Language Resources for K-12 Curriculum Development}
\author{Compiled by Claude Sonnet 5 (claude-sonnet-5) for Andr\'as Kornai}
\date{}
\begin{document}
\maketitle

\section{Introduction}
Akei (also Tasiriki; ISO 639-3 tsr; Glottocode akei1237) is an Oceanic language spoken on the southwest coast of Espiritu Santo, Vanuatu. Speaker figures conflict: about 650 (1981, via Wikipedia) and about 4,000 (Tryon 2010, via the Espiritu Santo languages page). Some sources treat it as a dialect of a proposed Southwest Santo language. The searches found no Akei-specific school materials.

\section{Primary Educational Resources}
No Akei primers or textbooks were found. For national context, Vanuatu has a policy of vernacular language use in early schooling:
\begin{itemize}
\item Franjieh et al. (2025), ``Vernacular language use in schools: Teachers' perspectives on implementing the Vanuatu national language policy''. It covers other languages, not Akei: \url{https://nzlingsoc.org/wp-content/uploads/2025/12/3-5-Franjieh-et-al-2025-Vernacular-language-use-in-schools.pdf}
\item ``Vernacular literacy development in Vanuatu: Comparative case studies''. It covers North Ambrym and Ske (Pentecost), not Akei, and is a model for community-led orthography and teacher training: \url{https://scholarspace.manoa.hawaii.edu/items/0e64353d-05e0-4cd1-8a5a-5acb252d9dc7}
\end{itemize}

\section{Linguistic and Reference Materials}
\begin{itemize}
\item Akei language (Wikipedia): \url{https://en.wikipedia.org/wiki/Akei_language}
\item Akei entry in the Tlopo/CLLD database: \url{https://tlopo.clld.org/languages/akei}
\item Espiritu Santo languages (Wikipedia): \url{https://en.wikipedia.org/wiki/Espiritu_Santo_languages}
\end{itemize}

\section{Religious and Community Resources}
Joshua Project profile of the Akei (Tasiriki) people: \url{https://joshuaproject.net/people_groups/10222}. No scripture or community organisation was found.

\section{Digital and Audio Resources}
No audio, video or software resources specific to Akei were found.

\section{Curriculum Development Context}
The reference description says Akei is used as a language of instruction in some schools. The searches did not confirm this. A project would need to find out from local schools and the Vanuatu Ministry of Education what exists, and should establish an orthography and basic literacy materials with the community.
\end{document}